\section{Related Work}

\paragraph{Scaling laws and spectral learning curves.}
Empirical neural scaling laws were documented by
\cite{kaplan2020scaling} and refined for compute-optimal training by
\cite{hoffmann2022training}. Spectral theories explain learning curves through
resolution of covariance or kernel eigenmodes
\cite{bordelon2020spectrum,bahri2024explaining}. Solvable random-feature
models connect power-law data structure to model- and data-scaling exponents
\cite{maloney2022solvable}, while feature learning can alter those exponents
outside the kernel regime \cite{bordelon2024feature}.

Our closest starting point is the infinite-dimensional sketched linear model
of \cite{lin2024scaling}, which proves a two-resource law for one-pass SGD on
i.i.d. data. \cite{lin2025reuse} adds a source exponent and shows how
multi-pass optimization can improve the iteration term when data are reused.
We instead keep the marginal task fixed and modify the temporal arrival
process. Reusing a persistent block cannot reveal a mode that has not arrived;
this creates an innovation-limited spectral cutoff.

\paragraph{Regression with dependent samples.}
Least squares under Markovian sampling is known to depend on mixing time, and
worst-case minimax rates can be attained by data dropping; structured chains
can benefit from experience replay \cite{nagaraj2020markovian}. Finite-time
SGD analyses and recent single-trajectory lower bounds likewise quantify
Markov dependence through global chain parameters
\cite{doan2020finite,sun2025single}. Recent high-dimensional ridge theory
allows correlated Gaussian examples and derives sharp risk and
cross-validation formulas \cite{atanasov2024correlated}. Its
matrix-Gaussian model has separable covariance
$\E[X_{ti}X_{sj}]=K_{ts}\Sigma_{ij}$: every feature direction inherits the
same sample-correlation matrix. Our construction is nonseparable. Mode $j$
has its own correlation time $\tau_j$, and the alignment between
$\lambda_j$ and $\tau_j$ is precisely what changes the exponent.

\paragraph{Infinite occupancy.}
At innovation times, the problem is an infinite-urn scheme with frequencies
$q_j$. Missing-mass and occupancy asymptotics under regular variation are
well studied \cite{gnedin2007occupancy,benhamou2017concentration}. Our risk is
a \emph{target-weighted} missing mass: unseen mode $j$ costs $s_j$, which may
have a different exponent from $q_j$. We use this weighted structure to derive
the model--data scaling law and minimax regression interpretation.

